% ===========================================================================
%  medimg — Compressed Sensing Reconstruction for Accelerated MRI and CT
%  Technical Note
%  Version 0.1.0 — July 2026
% ===========================================================================

\documentclass[11pt,a4paper]{article}

% --- Packages ----------------------------------------------------------------
\usepackage[american]{babel}
\usepackage[utf8]{inputenc}
\usepackage[T1]{fontenc}
\usepackage{amsmath,amssymb}
\usepackage{booktabs}
\usepackage[margin=2.5cm]{geometry}
\usepackage[numbers,sort&compress]{natbib}
\usepackage[colorlinks=true,linkcolor=black,citecolor=black,urlcolor=blue]{hyperref}
\usepackage{enumitem}
\usepackage{xcolor}
\usepackage{abstract}

% --- Title helpers -----------------------------------------------------------
\newcommand{\seprule}{\vspace{4pt}\hrule\vspace{4pt}}

% ===========================================================================
\begin{document}

% --- Title block -------------------------------------------------------------
\begin{center}
{\LARGE\bfseries medimg}\\[4pt]
{\large Compressed Sensing Reconstruction for Accelerated MRI and CT}\\[8pt]
{\large Abd'gafar Tunde Tiamiyu}\\[2pt]
School of Mathematics, Jilin University, Changchun 130012, China\\[2pt]
\texttt{abdgafartunde@yahoo.com}\\[8pt]
Version 0.1.0 — July 2026\\[2pt]
DOI: \texttt{10.5281/zenodo.21321200}
\seprule
\end{center}

% --- Summary -----------------------------------------------------------------
\begin{abstract}
\noindent\textbf{medimg} is a research-grade Python package for compressed
sensing and model-based image reconstruction in accelerated magnetic
resonance imaging (MRI) and sparse-view computed tomography (CT).  It
implements three classical variational methods — Tikhonov regularisation
(conjugate gradient), isotropic total variation (Chambolle-Pock
primal-dual), and wavelet-domain $\ell_1$ minimisation (FISTA) — with
clean operator-based abstractions that work across both modalities.
Variable-density Cartesian $k$-space under-sampling (Lustig et~al., 2007)
and sparse-angle parallel-beam Radon projections (Sidky \& Pan, 2008)
serve as the forward models.  The package includes two demonstration
scripts that generate publication-quality comparison figures, and a
suite of 41 unit tests.
\end{abstract}

% ===========================================================================
\section{Overview}
% ===========================================================================

Accelerated acquisition in MRI and CT reduces scan time and radiation
dose, but the resulting under-sampled data lead to an ill-posed linear
inverse problem.  In MRI, $k$-space is sub-sampled below the Nyquist
rate; in CT, the number of projection angles is reduced.  Compressed
sensing (CS) theory \cite{Candes2006,Donoho2006} guarantees exact
recovery under sparsity and incoherence assumptions, provided a suitable
nonlinear reconstruction algorithm is employed.

This technical note describes the forward models (Section~2), the
compressed sensing principles that underpin the reconstruction
(Section~3), the three variational methods implemented in the package
(Section~4), and the software structure and usage (Section~5).

% ===========================================================================
\section{Forward Models}
% ===========================================================================

\subsection{MRI}

The single-coil MRI measurement process is the 2-D Fourier transform
sampled at discrete $k$-space locations.  Let $\mathcal{F}$ denote the
unitary 2-D DFT, $M\in\{0,1\}^{n_r\times n_c}$ a binary under-sampling
mask, and $\eta\sim\mathcal{CN}(0,\sigma^2 I)$ additive complex Gaussian
noise.  The forward operator is
\begin{equation}\label{eq:mri_forward}
A_{\text{MRI}}(x)=M\odot\mathcal{F}(x),\qquad
y=A_{\text{MRI}}(x)+\eta,
\end{equation}
with adjoint
$A_{\text{MRI}}^*(y)=\operatorname{Re}\bigl(\mathcal{F}^{-1}(M\odot y)\bigr)$.

The variable-density mask \cite{Lustig2007} fully samples a central
calibration band of $n_c$ low-frequency phase-encode lines and randomly
sub-samples higher frequencies with probability
$P(k_y)\propto\bigl(1-(|k_y|/k_y^{\max})^2\bigr)^2$, exploiting the
energy concentration of natural images at low spatial frequencies.
The under-sampling ratio $R=\|M\|_0/(n_r n_c)$ controls the
acceleration factor $1/R$.

\subsection{CT}

For parallel-beam geometry, the forward operator is the 2-D Radon
transform $\mathcal{R}$.  With projection angles $\theta_j\in[0,\pi)$
and detector coordinate~$s$,
\begin{equation}\label{eq:ct_forward}
A_{\text{CT}}(x)(\theta_j,s)=\mathcal{R}x(\theta_j,s)
=\int_{-\infty}^{\infty}x(s\cos\theta_j-t\sin\theta_j,\;
s\sin\theta_j+t\cos\theta_j)\,dt.
\end{equation}
The adjoint $A_{\text{CT}}^*=\mathcal{R}^*$ is the unfiltered
back-projection (not the FBP reconstruction, which applies a ramp
filter).  Sparse-view CT uses $N_\theta\ll N_r$ projection angles,
making the problem under-determined.  We normalise the operator pair by
$1/\sqrt{\pi N_\theta}$ so that $\|A\|\approx 1$, matching the MRI
scaling and making regularisation parameters transferable across
modalities.

% ===========================================================================
\section{Compressed Sensing Principles}
% ===========================================================================

Compressed sensing \cite{Candes2006,Donoho2006} guarantees exact or
near-exact recovery of a sparse signal from under-sampled linear
measurements under two conditions:
\begin{enumerate}[label=(\arabic*)]
\item \textbf{Sparsity.}  The image $x$ admits a sparse representation
  in a transform domain~$\Psi$: $x=\Psi^*c$, where $c$ has few non-zero
  entries.
\item \textbf{Incoherence.}  The sensing basis (Fourier for MRI; pixel
  basis for CT) and the sparsity basis (wavelet; finite differences)
  are sufficiently incoherent.
\end{enumerate}
Under these conditions, the $\ell_1$-regularised problem
\begin{equation}\label{eq:cs_objective}
\min_x\;\frac12\|Ax-y\|_2^2+\lambda\|\Psi x\|_1
\end{equation}
recovers the true image with high probability when $M\gtrsim
s\log(N/s)$ measurements are acquired, where $s$ is the sparsity
level.

% ===========================================================================
\section{Variational Reconstruction Methods}
% ===========================================================================

\subsection{Tikhonov Regularisation ($\ell_2$)}

The simplest convex regulariser penalises the $\ell_2$ energy:
\begin{equation}\label{eq:tikhonov}
\min_x\;\frac12\|Ax-y\|_2^2+\frac{\lambda}{2}\|x\|_2^2.
\end{equation}
The solution satisfies the normal equation
$(A^*A+\lambda I)x=A^*y$, a symmetric positive-definite linear system
solved by the conjugate gradient method.  Tikhonov regularisation
produces smooth reconstructions; it does not promote sparsity or
preserve edges.

\subsection{Total Variation (TV)}

The discrete isotropic total variation is
\begin{equation}\label{eq:tv}
\operatorname{TV}(x)=\|\nabla x\|_{2,1}
=\sum_{i,j}\sqrt{(\nabla_1 x)_{i,j}^2+(\nabla_2 x)_{i,j}^2},
\end{equation}
where $\nabla_1,\nabla_2$ are forward finite-difference operators.
TV penalises the $\ell_1$ norm of the gradient magnitude, promoting
piecewise-constant reconstructions with sharp edges.

The reconstruction problem
$\min_x\frac12\|Ax-y\|_2^2+\alpha\operatorname{TV}(x)$
is solved via the Chambolle-Pock primal-dual algorithm
\cite{Chambolle2011}.  Defining the saddle-point formulation
\begin{equation*}
\min_x\max_{p:\,\|p\|_{2,\infty}\le\alpha}
\frac12\|Ax-y\|_2^2+\langle\nabla x,p\rangle,
\end{equation*}
the algorithm alternates dual ascent (pointwise $\ell_2$-ball
projection), primal descent with the adjoint of $[A;\;\nabla]$, and
over-relaxation.  The step sizes $\tau=\sigma=0.99/\sqrt{\|A\|^2+8}$
guarantee convergence via $\tau\sigma\|K\|^2<1$.

\subsection{Wavelet-Domain $\ell_1$ Minimisation (FISTA)}

Let $W$ be a sparsifying 2-D discrete wavelet transform (e.g.,
Haar) and $W^*$ its adjoint.  The CS objective~\eqref{eq:cs_objective}
with $\Psi=W$ is a composite convex problem solved by FISTA
\cite{Beck2009} with optimal $O(1/k^2)$ convergence:
\begin{align}
\tilde{x}^k &= z^k - \tfrac{1}{L}A^*(Az^k-y), \label{eq:fista_grad}\\
x^{k+1} &= W^*\bigl(S_{\lambda/L}(W\tilde{x}^k)\bigr), \label{eq:fista_prox}\\
t_{k+1} &= \tfrac12\bigl(1+\sqrt{1+4t_k^2}\bigr),\qquad
z^{k+1}=x^{k+1}+\tfrac{t_k-1}{t_{k+1}}(x^{k+1}-x^k). \label{eq:fista_mom}
\end{align}
Here $L=\|A\|^2$ (estimated by power iteration when not known
analytically) and
$S_\tau(c)=\operatorname{sign}(c)\max(|c|-\tau,0)$ is element-wise
soft-thresholding.

The Haar wavelet is used as the default sparsifying transform because
its piecewise-constant basis functions match the structure of the
standard Shepp-Logan phantom, yielding substantially sparser
coefficient vectors than smoother wavelets (e.g., Daubechies-4).

% ===========================================================================
\section{Package Structure and Availability}
% ===========================================================================

The package is organised into flat modules:

\begin{center}
\begin{tabular}{@{}ll@{}}
\toprule
\textbf{Module} & \textbf{Description} \\
\midrule
\texttt{medimg.forward}    & MRI (Fourier + mask) and CT (Radon) operators \\
\texttt{medimg.reconstruct}& Tikhonov (CG), TV (Chambolle-Pock), CS (FISTA) \\
\texttt{medimg.transforms} & Wavelet analysis/synthesis, soft-thresholding \\
\texttt{medimg.phantoms}   & Shepp-Logan, custom circle phantoms \\
\texttt{medimg.metrics}    & PSNR, SSIM (Wang et~al., 2004) \\
\texttt{medimg.plotting}   & Side-by-side comparisons, error maps \\
\bottomrule
\end{tabular}
\end{center}

\subsection{Installation}

\begin{verbatim}
git clone https://github.com/abdgafartunde/medical-imaging
cd medical-imaging
pip install -e ".[dev]"
\end{verbatim}

Dependencies: NumPy ($\ge1.24$), SciPy ($\ge1.10$), Matplotlib
($\ge3.7$), scikit-image ($\ge0.20$), PyWavelets ($\ge1.4$).

\subsection{Quick Start}

\begin{verbatim}
from medimg.phantoms import shepp_logan
from medimg.forward import generate_cartesian_mask, mri_forward, mri_adjoint
from medimg.reconstruct import cs_fista
from medimg.transforms import wavelet_forward, wavelet_adjoint

phantom = shepp_logan((256, 256))
mask    = generate_cartesian_mask(phantom.shape, sampling_rate=0.25,
                                  calib_lines=24, rng=42)
A  = lambda x: mri_forward(x, mask)
At = lambda y: mri_adjoint(y, mask)
y  = mri_forward(phantom, mask, noise_std=0.001, rng=42)

W  = lambda x: wavelet_forward(x, wavelet='haar', level=4)
Wt = lambda c: wavelet_adjoint(c, wavelet='haar')
x_cs = cs_fista(A, At, y, W, Wt, lambd=1e-4, shape=phantom.shape)
\end{verbatim}

\subsection{Examples and Tests}

\begin{verbatim}
python examples/01_mri_cs.py        # CS-MRI demo
python examples/02_ct_sparse_view.py # sparse-view CT demo
pytest tests/ -v                     # 41 unit tests
\end{verbatim}

% ===========================================================================
\section{Image Quality Metrics}
% ===========================================================================

\textbf{PSNR} (Peak Signal-to-Noise Ratio, in dB):
\begin{equation*}
\operatorname{PSNR}(x,x_{\text{ref}})=
20\log_{10}(\operatorname{range})-
10\log_{10}\!\Bigl(\tfrac{1}{N}\|x-x_{\text{ref}}\|_2^2\Bigr).
\end{equation*}

\textbf{SSIM} (Structural Similarity) \cite{Wang2004}:
\begin{equation*}
\operatorname{SSIM}(x,x_{\text{ref}})=
\frac{(2\mu_x\mu_{\text{ref}}+C_1)(2\sigma_{x,\text{ref}}+C_2)}
{(\mu_x^2+\mu_{\text{ref}}^2+C_1)(\sigma_x^2+\sigma_{\text{ref}}^2+C_2)},
\end{equation*}
where $\mu,\sigma$ are local means and standard deviations over
$11\times11$ Gaussian windows, and $C_1,C_2$ are stabilisation
constants.

% ===========================================================================
\begin{thebibliography}{10}

\bibitem{Lustig2007}
M.~Lustig, D.~Donoho, and J.~M.~Pauly.
\newblock Sparse MRI: the application of compressed sensing for rapid MR
  imaging.
\newblock \emph{Magnetic Resonance in Medicine}, 58(6):1182--1195, 2007.

\bibitem{Sidky2008}
E.~Y.~Sidky and X.~Pan.
\newblock Image reconstruction in circular cone-beam computed tomography by
  constrained, total-variation minimization.
\newblock \emph{Physics in Medicine and Biology}, 53(17):4777--4807, 2008.

\bibitem{Candes2006}
E.~J.~Cand\`es, J.~Romberg, and T.~Tao.
\newblock Robust uncertainty principles: exact signal reconstruction from
  highly incomplete frequency information.
\newblock \emph{IEEE Transactions on Information Theory}, 52(2):489--509,
  2006.

\bibitem{Donoho2006}
D.~L.~Donoho.
\newblock Compressed sensing.
\newblock \emph{IEEE Transactions on Information Theory}, 52(4):1289--1306,
  2006.

\bibitem{Chambolle2011}
A.~Chambolle and T.~Pock.
\newblock A first-order primal-dual algorithm for convex problems with
  applications to imaging.
\newblock \emph{Journal of Mathematical Imaging and Vision},
  40(1):120--145, 2011.

\bibitem{Beck2009}
A.~Beck and M.~Teboulle.
\newblock A fast iterative shrinkage-thresholding algorithm for linear
  inverse problems.
\newblock \emph{SIAM Journal on Imaging Sciences}, 2(1):183--202, 2009.

\bibitem{Wang2004}
Z.~Wang, A.~C.~Bovik, H.~R.~Sheikh, and E.~P.~Simoncelli.
\newblock Image quality assessment: from error visibility to structural
  similarity.
\newblock \emph{IEEE Transactions on Image Processing}, 13(4):600--612,
  2004.

\end{thebibliography}

\end{document}
